% questions.tex —— 2014 年·江苏卷 · 题目内容片段（无答案）
% 由 tools/generate_projects.py 依来源总账自 2014/2014_数学真题_江苏_试题.md 生成，勿手改（改源 MD 后重跑）


\begin{gknotice}
\textbf{来源说明}：本题面依据固定源 PDF 的两页 A3 题面逐题转录，原卷印刷页码为 889—890。第 889 面为填空题 1—14 与解答题 15—18，第 890 面为解答题 19—20 与附加题 21—23（第 21 题为四选二）。

\textbf{图形说明}：第 3 题程序框图、第 6 题频率分布直方图、第 12 题平行四边形、第 16 题三棱锥、第 17 题椭圆、第 18 题河桥与圆形保护区、第 21 题【A】圆，均由项目内 TikZ 矢量源重绘；其余各题原页无图。

\textbf{解析归属}：固定源不附答案或解析；配套答案与解析均由本库按本题面独立复算。
\end{gknotice}

\gkpart{一、填空题}

\begin{question}
\gkti{1}已知集合 $A=\{-2,-1,3,4\}$，$B=\{-1,2,3\}$，则 $A\cap B=$\gkblank{}。

\end{question}
\begin{question}
\gkti{2}已知复数 $z=(5+2\mathrm i)^2$（$\mathrm i$ 为虚数单位），则 $z$ 的实部为\gkblank{}。

\Needspace{0.4\textheight}
\end{question}
\begin{question}
\gkti{3}如图是一个算法流程图，则输出的 $n$ 的值是\gkblank{}。
\nopagebreak[4]
\begin{center}% 2014 年江苏卷填空题第 3 题：算法流程图，输出 n 的值。
% 来源 PDF：普通高考/2014/2014江苏.pdf，印面 889。
\begin{tikzpicture}[>=Stealth, node distance=7mm, font=\small,
  term/.style={draw, rounded corners=4pt, minimum width=20mm, minimum height=8mm},
  proc/.style={draw, minimum width=27mm, minimum height=8mm},
  dec/.style={draw, diamond, aspect=2.2, minimum width=32mm, minimum height=14mm,
    inner sep=1pt},
  io/.style={draw, trapezium, trapezium left angle=70, trapezium right angle=110,
    minimum width=27mm, minimum height=8mm, inner sep=2pt}]
  \node[term] (start) {开始};
  \node[proc, below=of start] (init) {$n\leftarrow 0$};
  \node[proc, below=of init] (inc) {$n\leftarrow n+1$};
  \node[dec, below=of inc] (test) {$2^n>20$};
  \node[io, below=of test] (out) {输出 $n$};
  \node[term, below=of out] (stop) {结束};

  \draw[->] (start) -- (init);
  \draw[->] (init) -- (inc);
  \draw[->] (inc) -- (test);
  \draw[->] (test) -- node[right, inner sep=1pt] {Y} (out);
  \draw[->] (out) -- (stop);
  \draw[->] (test.east) -- ++(13mm,0) |- node[pos=0.25, right, inner sep=1pt] {N}
    (inc.north);
\end{tikzpicture}
\end{center}
\nopagebreak[4]
\end{question}
\begin{question}
\gkti{4}从 $1,2,3,6$ 这 4 个数中一次随机地取 2 个数，则所取 2 个数的乘积为 6 的概率是\gkblank{}。

\end{question}
\begin{question}
\gkti{5}已知函数 $y=\cos x$ 与 $y=\sin(2x+\varphi)$（$0\leqslant\varphi<\pi$），它们的图象有一个横坐标为 $\dfrac\pi3$ 的交点，则 $\varphi$ 的值是\gkblank{}。

\Needspace{0.4\textheight}
\end{question}
\begin{question}
\gkti{6}为了了解一片经济林的生长情况，随机抽测了其中 60 株树木的底部周长（单位：cm），所得数据均在区间 $[80,130]$ 上，其频率分布直方图如图所示，则在抽测的 60 株树木中，有\gkblank{}株树木的底部周长小于 100 cm。
\nopagebreak[4]
\begin{center}% 2014 年江苏卷填空题第 6 题：60 株树木底部周长的频率分布直方图。
% 来源 PDF：普通高考/2014/2014江苏.pdf，印面 889。组距 10 cm，五组高度依次为
% 0.015、0.025、0.030、0.020、0.010。
\begin{tikzpicture}[>=Stealth, semithick, font=\small, scale=0.85, transform shape]
  % 纵轴按 180 cm / 单位高度绘制
  \draw[->] (0,0) -- (0,6.5) node[above] {$\dfrac{\text{频率}}{\text{组距}}$};
  \draw[->] (0,0) -- (10.2,0);
  \node[below right, inner sep=4pt] at (10.0,-0.15) {底部周长/cm};
  % 纵轴刻度与虚线
  \foreach \h/\lab in {1.8/0.010, 2.7/0.015, 3.6/0.020, 4.5/0.025, 5.4/0.030}{
    \draw[dashed] (0,\h) -- (9.6,\h);
    \node[left, inner sep=3pt] at (0,\h) {\lab};
  }
  % 五根矩形：组距 10 cm 对应宽度 1.6 cm
  \draw (1.6,0) rectangle (3.2,2.7);
  \draw (3.2,0) rectangle (4.8,4.5);
  \draw (4.8,0) rectangle (6.4,5.4);
  \draw (6.4,0) rectangle (8.0,3.6);
  \draw (8.0,0) rectangle (9.6,1.8);
  % 横轴刻度
  \node[below, inner sep=3pt] at (0.15,0) {$0$};
  \foreach \x/\lab in {1.6/80, 3.2/90, 4.8/100, 6.4/110, 8.0/120, 9.6/130}{
    \node[below, inner sep=3pt] at (\x,0) {\lab};
  }
  % 断裂符
  \draw (0.55,0.28) -- (0.95,-0.28);
  \draw (0.85,0.28) -- (1.25,-0.28);
\end{tikzpicture}
\end{center}
\nopagebreak[4]
\end{question}
\begin{question}
\gkti{7}在各项均为正数的等比数列 $\{a_n\}$ 中，若 $a_2=1$，$a_8=a_6+2a_4$，则 $a_6$ 的值是\gkblank{}。

\end{question}
\begin{question}
\gkti{8}设甲、乙两个圆柱的底面积分别为 $S_1$、$S_2$，体积分别为 $V_1$、$V_2$，若它们的侧面积相等，且 $\dfrac{S_1}{S_2}=\dfrac94$，则 $\dfrac{V_1}{V_2}$ 的值是\gkblank{}。

\end{question}
\begin{question}
\gkti{9}在平面直角坐标系 $xOy$ 中，直线 $x+2y-3=0$ 被圆 $(x-2)^2+(y+1)^2=4$ 截得的弦长为\gkblank{}。

\end{question}
\begin{question}
\gkti{10}已知函数 $f(x)=x^2+mx-1$，若对于任意 $x\in[m,m+1]$，都有 $f(x)<0$ 成立，则实数 $m$ 的取值范围是\gkblank{}。

\end{question}
\begin{question}
\gkti{11}在平面直角坐标系 $xOy$ 中，若曲线 $y=ax^2+\dfrac bx$（$a$，$b$ 为常数）过点 $P(2,-5)$，且该曲线在点 $P$ 处的切线与直线 $7x+2y+3=0$ 平行，则 $a+b$ 的值是\gkblank{}。

\Needspace{0.4\textheight}
\end{question}
\begin{question}
\gkti{12}如图，在平行四边形 $ABCD$ 中，已知 $AB=8$，$AD=5$，$\overrightarrow{CP}=3\overrightarrow{PD}$，$\overrightarrow{AP}\cdot\overrightarrow{BP}=2$，则 $\overrightarrow{AB}\cdot\overrightarrow{AD}$ 的值是\gkblank{}。
\nopagebreak[4]
\begin{center}% 2014 年江苏卷填空题第 12 题：平行四边形 ABCD 中 P 在 DC 上，
% CP=3PD，连 AP、BP。来源 PDF：普通高考/2014/2014江苏.pdf，印面 889。
\begin{tikzpicture}[>=Stealth, semithick, font=\small, scale=0.95]
  \coordinate (A) at (0,0);
  \coordinate (B) at (6,0);
  \coordinate (C) at (7.8,3.4);
  \coordinate (D) at (1.8,3.4);
  \coordinate (P) at ($(D)!0.25!(C)$);
  \draw (A) -- (B) -- (C) -- (D) -- cycle;
  \draw (A) -- (P) -- (B);
  \fill (A) circle (1.2pt) node[below left] {$A$};
  \fill (B) circle (1.2pt) node[below right] {$B$};
  \fill (C) circle (1.2pt) node[right] {$C$};
  \fill (D) circle (1.2pt) node[left] {$D$};
  \fill (P) circle (1.2pt) node[above] {$P$};
\end{tikzpicture}
\end{center}
\nopagebreak[4]
\end{question}
\begin{question}
\gkti{13}已知 $f(x)$ 是定义在 $\mathbf R$ 上且周期为 3 的函数，当 $x\in[0,3)$ 时，$f(x)=\left|x^2-2x+\dfrac12\right|$。若函数 $y=f(x)-a$ 在区间 $[-3,4]$ 上有 10 个零点（互不相同），则实数 $a$ 的取值范围是\gkblank{}。

\end{question}
\begin{question}
\gkti{14}若 $\triangle ABC$ 的内角满足 $\sin A+\sqrt2\sin B=2\sin C$，则 $\cos C$ 的最小值是\gkblank{}。

\end{question}
\gkpart{二、解答题}

\begin{question}
\gkti{15}已知 $\alpha\in\left(\dfrac\pi2,\pi\right)$，$\sin\alpha=\dfrac{\sqrt5}{5}$。

（1） 求 $\sin\left(\dfrac\pi4+\alpha\right)$ 的值；

（2） 求 $\cos\left(\dfrac{5\pi}6-2\alpha\right)$ 的值。

\Needspace{0.4\textheight}
\end{question}
\begin{question}
\gkti{16}如图，在三棱锥 $P-ABC$ 中，$D$，$E$，$F$ 分别为棱 $PC$，$AC$，$AB$ 的中点。已知 $PA\perp AC$，$PA=6$，$BC=8$，$DF=5$。

（1） 求证：直线 $PA\parallel$ 平面 $DEF$；

（2） 求证：平面 $BDE\perp$ 平面 $ABC$。
\nopagebreak[4]
\begin{center}% 2014 年江苏卷解答题第 16 题：三棱锥 P-ABC，D、E、F 分别为 PC、AC、AB 中点。
% PA⊥AC，PA=6，BC=8，DF=5。来源 PDF：普通高考/2014/2014江苏.pdf，印面 889。
\begin{tikzpicture}[>=Stealth, semithick, font=\small, scale=0.95]
  \coordinate (P) at (0,4.6);
  \coordinate (A) at (0,1.5);
  \coordinate (B) at (1.3,0);
  \coordinate (C) at (4.6,1.7);
  \coordinate (D) at ($(P)!0.5!(C)$);
  \coordinate (E) at ($(A)!0.5!(C)$);
  \coordinate (F) at ($(A)!0.5!(B)$);
  % 棱
  \draw (P) -- (A);
  \draw (P) -- (B);
  \draw (P) -- (C);
  \draw (A) -- (B);
  \draw (B) -- (C);
  \draw[dashed] (A) -- (C);
  % 辅助线
  \draw[dashed] (D) -- (E);
  \draw[dashed] (E) -- (F);
  \draw[dashed] (F) -- (D);
  \draw[dashed] (B) -- (E);
  \draw[dashed] (B) -- (D);
  % 顶点与中点标注
  \fill (P) circle (1.1pt) node[above] {$P$};
  \fill (A) circle (1.1pt) node[left] {$A$};
  \fill (B) circle (1.1pt) node[below] {$B$};
  \fill (C) circle (1.1pt) node[right] {$C$};
  \fill (D) circle (1.1pt) node[above right] {$D$};
  \fill (E) circle (1.1pt) node[below right] {$E$};
  \fill (F) circle (1.1pt) node[below left] {$F$};
\end{tikzpicture}
\end{center}
\nopagebreak[4]
\Needspace{0.4\textheight}
\end{question}
\begin{question}
\gkti{17}如图，在平面直角坐标系 $xOy$ 中，$F_1$，$F_2$ 分别是椭圆 $\dfrac{x^2}{a^2}+\dfrac{y^2}{b^2}=1$（$a>b>0$）的左、右焦点，顶点 $B$ 的坐标为 $(0,b)$，连接 $BF_2$ 并延长交椭圆于点 $A$，过点 $A$ 作 $x$ 轴的垂线交椭圆于另一点 $C$，连接 $F_1C$。

（1） 若点 $C$ 的坐标为 $\left(\dfrac43,\dfrac13\right)$，且 $BF_2=\sqrt2$，求椭圆的方程；

（2） 若 $F_1C\perp AB$，求椭圆离心率 $e$ 的值。
\nopagebreak[4]
\begin{center}% 2014 年江苏卷解答题第 17 题：椭圆 x^2/a^2+y^2/b^2=1，B(0,b)，BF2 延长交椭圆于 A，
% 过 A 作 x 轴垂线交椭圆于 C，连 F1C。示意图取 a:b:c=2.8:2.2:1.73（非按离心率作答值定图）。
% 来源 PDF：普通高考/2014/2014江苏.pdf，印面 889。
\begin{tikzpicture}[>=Stealth, semithick, font=\small, scale=1.0]
  % 坐标轴
  \draw[->] (-3.6,0) -- (3.75,0) node[below, inner sep=2pt] {$x$};
  \draw[->] (0,-2.75) -- (0,2.95);
  % 椭圆
  \draw[thick] (0,0) ellipse (2.8 and 2.2);
  % 点：A、C 为过 A 的竖直线与椭圆的交点
  \coordinate (O) at (0,0);
  \coordinate (B) at (0,2.2);
  \coordinate (F1) at (-1.73,0);
  \coordinate (F2) at (1.73,0);
  \coordinate (A) at (2.506,-0.982);
  \coordinate (C) at (2.506,0.982);
  % 连线
  \draw (B) -- (A);
  \draw (F1) -- (C);
  \draw (A) -- (C);
  % 标注
  \fill (O) circle (1pt) node[below left, inner sep=2pt] {$O$};
  \fill (B) circle (1pt) node[above] {$B$};
  \fill (F1) circle (1pt) node[below, inner sep=3pt] {$F_1$};
  \fill (F2) circle (1pt) node[below, inner sep=3pt] {$F_2$};
  \fill (A) circle (1pt) node[below right, inner sep=2pt] {$A$};
  \fill (C) circle (1pt) node[above right, inner sep=2pt] {$C$};
\end{tikzpicture}
\end{center}
\nopagebreak[4]
\Needspace{0.4\textheight}
\end{question}
\begin{question}
\gkti{18}如图，为了保护河上古桥 $OA$，规划建一座新桥 $BC$，同时设立一个圆形保护区。规划要求：新桥 $BC$ 与河岸 $AB$ 垂直；保护区的边界为圆心 $M$ 在线段 $OA$ 上并与 $BC$ 相切的圆，且古桥两端 $O$ 和 $A$ 到该圆上任意一点的距离均不少于 80 m。经测量，点 $A$ 位于点 $O$ 正北方向 60 m 处，点 $C$ 位于点 $O$ 正东方向 170 m 处（$OC$ 为河岸），$\tan\angle BCO=\dfrac43$。

（1） 求新桥 $BC$ 的长；

（2） 当 $OM$ 多长时，圆形保护区的面积最大？
\nopagebreak[4]
\begin{center}% 2014 年江苏卷解答题第 18 题：古桥 OA、新桥 BC、圆形保护区（圆心 M 在 OA 上，
% 与 BC 相切）。OA=60 m，OC=170 m，BC⊥AB。来源 PDF：普通高考/2014/2014江苏.pdf，印面 889。
% 示意比例：1 单位≈30 m，故 OA=2，OC≈5.67，B=(2.67,4)。
\begin{tikzpicture}[>=Stealth, semithick, font=\small, scale=1.0]
  \coordinate (O) at (0,0);
  \coordinate (A) at (0,2);
  \coordinate (M) at (0,1.45);
  \coordinate (C) at (5.67,0);
  \coordinate (B) at (2.67,4);
  % 保护区（虚线圆）
  \draw[dashed] (M) circle (3.66);
  % 河岸与新桥
  \draw (A) -- (B) -- (C);
  \draw (O) -- (C);
  \draw (O) -- (A);
  % 直角标记
  \draw ($(B)!0.55mm!(A)$) -- ($(B)!0.55mm!(A)!0.55mm!90:(B)$)
    -- ($(B)!0.55mm!(C)$);
  % 方向标
  \draw[->] (0,0) -- (0,4.72) node[above] {北};
  \draw[->] (0,0) -- (6.7,0) node[below] {东};
  % 60 m 标注
  \draw[<->] (-0.55,0) -- (-0.55,2);
  \node[left] at (-0.62,1.0) {$60\,\mathrm m$};
  % 170 m 标注
  \draw[<->] (0,-0.55) -- (5.67,-0.55);
  \node[below] at (2.84,-0.95) {$170\,\mathrm m$};
  % 点标注
  \fill (O) circle (1pt) node[below left, inner sep=2pt] {$O$};
  \fill (A) circle (1pt) node[left, inner sep=2pt] {$A$};
  \fill (M) circle (1.1pt) node[right, inner sep=2pt] {$M$};
  \fill (B) circle (1pt) node[above right] {$B$};
  \fill (C) circle (1pt) node[below right] {$C$};
\end{tikzpicture}
\end{center}
\nopagebreak[4]
\end{question}
\begin{question}
\gkti{19}已知函数 $f(x)=\mathrm e^x+\mathrm e^{-x}$，其中 $\mathrm e$ 是自然对数的底数。

（1） 证明：$f(x)$ 是 $\mathbf R$ 上的偶函数；

（2） 若关于 $x$ 的不等式 $mf(x)\leqslant \mathrm e^{-x}+m-1$ 在 $(0,+\infty)$ 上恒成立，求实数 $m$ 的取值范围；

（3） 已知正数 $a$ 满足：存在 $x_0\in[1,+\infty)$，使得 $f(x_0)<a\left(-x_0^3+3x_0\right)$ 成立。试比较 $\mathrm e^{a-1}$ 与 $a^{\mathrm e-1}$ 的大小，并证明你的结论。

\end{question}
\begin{question}
\gkti{20}设数列 $\{a_n\}$ 的前 $n$ 项和为 $S_n$。若对任意正整数 $n$，总存在正整数 $m$，使得 $S_n=a_m$，则称 $\{a_n\}$ 是“H 数列”。

（1） 若数列 $\{a_n\}$ 的前 $n$ 项和 $S_n=2^n$（$n\in\mathbf N^*$），证明：$\{a_n\}$ 是“H 数列”；

（2） 设 $\{a_n\}$ 是等差数列，其首项 $a_1=1$，公差 $d<0$。若 $\{a_n\}$ 是“H 数列”，求 $d$ 的值；

（3） 证明：对任意的等差数列 $\{a_n\}$，总存在两个“H 数列”$\{b_n\}$ 和 $\{c_n\}$，使得 $a_n=b_n+c_n$（$n\in\mathbf N^*$）成立。

\end{question}
\gkpart{三、附加题}

\Needspace{0.4\textheight}
\begin{question}
\gkti{21}四选二。

（1） 【A】如图，$AB$ 是圆 $O$ 的直径，$C$、$D$ 是圆 $O$ 上位于 $AB$ 异侧的两点。证明：$\angle OCB=\angle D$。
\nopagebreak[4]
\begin{center}% 2014 年江苏卷附加题 21A：AB 是圆 O 的直径，C、D 是圆 O 上位于 AB 异侧的两点。
% 来源 PDF：普通高考/2014/2014江苏.pdf，印面 890。
\begin{tikzpicture}[>=Stealth, semithick, font=\small, scale=1.0]
  \coordinate (O) at (0,0);
  \coordinate (A) at (140:2.1);
  \coordinate (B) at (-40:2.1);
  \coordinate (C) at (215:2.1);
  \coordinate (D) at (75:2.1);
  \draw[thick] (O) circle (2.1);
  \draw (A) -- (B);
  \draw (C) -- (D);
  \draw (A) -- (C);
  \draw (A) -- (D);
  \draw (B) -- (C);
  \draw (B) -- (D);
  \fill (O) circle (1pt) node[above right, inner sep=1pt] {$O$};
  \fill (A) circle (1.1pt) node[above left] {$A$};
  \fill (B) circle (1.1pt) node[below right] {$B$};
  \fill (C) circle (1.1pt) node[below left] {$C$};
  \fill (D) circle (1.1pt) node[above] {$D$};
\end{tikzpicture}
\end{center}
\nopagebreak[4]
（2） 【B】已知矩阵 $A=\begin{bmatrix}-1&2\\1&x\end{bmatrix}$，$B=\begin{bmatrix}1&1\\2&-1\end{bmatrix}$，向量 $\boldsymbol\alpha=\begin{bmatrix}2\\y\end{bmatrix}$，$x$，$y$ 为实数，若 $A\boldsymbol\alpha=B\boldsymbol\alpha$，求 $x+y$ 的值。

（3） 【C】在平面直角坐标系 $xOy$ 中，已知直线 $l$ 的参数方程为 $\begin{cases}x=1-\dfrac{\sqrt2}{2}t,\\y=2+\dfrac{\sqrt2}{2}t\end{cases}$（$t$ 为参数），直线 $l$ 与抛物线 $y^2=4x$ 交于 $A$，$B$ 两点，求线段 $AB$ 的长。

（4） 【D】已知 $x>0$，$y>0$，证明：$(1+x+y^2)(1+x^2+y)\geqslant 9xy$。

\end{question}
\begin{question}
\gkti{22}盒中共有 9 个球，其中有 4 个红球、3 个黄球和 2 个绿球，这些球除颜色外完全相同。

（1） 从盒中一次随机取出 2 个球，求取出的 2 个球颜色相同的概率 $P$；

（2） 从盒中一次随机取出 4 个球，其中红球、黄球、绿球的个数分别记为 $x_1$，$x_2$，$x_3$，随机变量 $X$ 表示 $x_1$，$x_2$，$x_3$ 中的最大数，求 $X$ 的概率分布和数学期望 $E(X)$。

\end{question}
\begin{question}
\gkti{23}已知函数 $f_0(x)=\dfrac{\sin x}{x}$（$x>0$），设 $f_n(x)$ 为 $f_{n-1}(x)$ 的导数，$n\in\mathbf N^*$。

（1） 求 $2f_1\left(\dfrac\pi2\right)+\dfrac\pi2f_2\left(\dfrac\pi2\right)$ 的值；

（2） 证明：对任意的 $n\in\mathbf N^*$，等式 $\left|nf_{n-1}\left(\dfrac\pi4\right)+\dfrac\pi4f_n\left(\dfrac\pi4\right)\right|=\dfrac{\sqrt2}{2}$ 都成立。

\end{question}